\subsection{Integral class five and an all-class relative lift}

We continue the finite-action test attached to \S52.2,
folio 744, scan 193. The next two results address different
lifting obligations. A.25 retains the particular class-two
conjugators of A.23 and the entire class-four family of A.24.
A.26 lifts the specified metabelian involution of A.22 through
the whole relative pro-$5$ completion, without prescribing
which of the earlier nonmetabelian conjugator pairs it extends.

\begin{statement}{A.25}{Proved}{The complete integral class-five family}
Let $N_5$ be the free class-five pro-$5$ quotient of the
index-three subgroup $K$ of A.23, and put
$Q_5^{[5]}=F/\ker(K\to N_5)$.
Consider all pairs of conjugators that extend the fixed
A.23 class-two pair, have the specified A.22 metabelian
images, and satisfy the two norm equations and the third
peripheral equation. Their complete parameter space is
an affine $\mathbb Z_5$-space of dimension $90$.
The resulting continuous automorphisms and their outer
classes have exact order two.

The calculation uses actual integral Hall group coordinates.
It retains all $20$ preceding parameters and introduces
$408$ central coordinates, giving an integral
$1428\times428$ affine system of rank $338$ over both
$\mathbb Q$ and $\mathbb F_5$.
The projection to the $20$-dimensional class-four family
has affine image of dimension $2$ and fibres of dimension
$88$. The class-four point whose $20$ parameters are zero
does not extend. The complete image, full affine kernel,
and all equations are recorded in the exact certificates.
\end{statement}
\begin{proof}
Keep the Schreier generators $a,b,c_0,d$ and their logarithmic
letters $A,B,C,D$ from A.23. Work in the rational tensor
algebra truncated after degree five. Its Lyndon Lie basis,
with $A<B<C<D$, has degree ranks
\[
                 4,\quad6,\quad20,\quad60,\quad204.
                                                               \tag{A--VII}
\]
For a letter $i$ let $g_i=\exp i$. For a longer Lyndon word
$w=uv$, with $v$ its longest proper Lyndon suffix, define
$g_w=[g_u,g_v]$, using $[g,h]=ghg^{-1}h^{-1}$.
The leading term of $\log g_w$ is the standard Lyndon
bracket $[w]$. The ordered group coordinates are
\[
       \prod_{|w|\leq5}g_w^{t_w},\qquad t_w\in\mathbb Z_5,
                                                               \tag{A--VIII}
\]
with degree first and lexicographic order within each degree.
These are coordinates in the free nilpotent pro-$5$ group,
not the assertion that the bare Lie lattice is BCH-closed.
Indeed, collection in each lower-central quotient identifies
it with the free integral Lie component. Changing its basis
to the Lyndon brackets is an integral basis change; lifting
those basis elements by the displayed group commutators
gives (A--VIII) inductively. Passing from integer coordinates
to their $5$-adic limits gives the free nilpotent pro-$5$
completion. The Hall normal form and this pro-$p$
interpretation are also stated in
\citet[\S4, displayed normal form and Remark 4.1]{A_CasalsHall2018}.
No naive reduction of all Hall coordinates modulo $5^n$
is used as a group quotient here.

The program checks all $294$ leading terms and all the
coefficients of the individual $\log g_w$; those coefficients
are $5$-integral. This does not imply that logarithms of
arbitrary products have integral Lie coordinates.
In particular a product of degree-one exponentials may
have degree-five logarithmic denominators divisible by $5$.

Let $(u_4,v_4)$ be the class-four particular solution in
the A.24 certificate. Collect $\exp u_4,\exp v_4$ through
degree four, and take the products (A--VIII) with these
first $90$ coordinates and zero final $204$ coordinates.
Call the resulting actual group elements $f_0,h_0$, with
rational logarithms $u_0,v_0$.
Let $(\delta u_i,\delta v_i)$, $0\leq i<20$, be the
entire A.24 nullspace basis. Every such correction has
degree four. All possible class-five pairs above that
family consequently have logarithms
\[
\begin{aligned}
 u={}&u_0+\sum_{i=0}^{19}q_i
       \bigl(\delta u_i+\tfrac12[u_1,\delta u_i]\bigr)+\alpha,\\
 v={}&v_0+\sum_{i=0}^{19}q_i
       \bigl(\delta v_i+\tfrac12[v_1,\delta v_i]\bigr)+\beta,
\end{aligned}                                                   \tag{A--IX}
\]
where $q_i\in\mathbb Z_5$, $\alpha,\beta\in L_5(\mathbb Z_5)$,
$u_1=\frac23(A+B-D)$, and $v_1=\frac23(A+C-D)$.
This parameterization is integral because it is exactly
$f_0\exp(\sum q_i\delta u_i)\exp\alpha$ and the analogous
product for $h_0$. The group of elements of weights at
least four is abelian here, and its integral logarithm
is $L_4\oplus L_5$, as follows from the unitriangular
Hall logarithms. Conversely, after the class-four
restriction is specified, every remaining choice differs
by an arbitrary central element of $L_5(\mathbb Z_5)$.
Thus none of the $20$ earlier parameters has been fixed.

The second-derived condition also requires an integral
comparison. In the rational Lie algebra put
\[
 H_{\mathbb Q}=\langle B-C,\ C-D+A\rangle_{\mathbb Q}
                      +\bigoplus_{j=2}^5L_j(\mathbb Q),
 \qquad H_{2,\mathbb Q}=[H_{\mathbb Q},H_{\mathbb Q}].
\]
The bracket ideal has rank $252$. The integral bracket
matrix has the same rank modulo $5$, and its echelon
basis is $5$-integral, so its span is the saturated lattice
$H_{2,\mathbb Q}\cap L_{\geq2}(\mathbb Z_5)$.
To identify this with $\log((Q_5^{[5]})'')$, observe first
that $N_5'$ has nilpotency class at most two. Its Hall
logarithms of weights at least two form a $5$-integral
unitriangular basis, and the only BCH denominator needed
inside this subgroup is $2$. Hence
$\log N_5'=L_{\geq2}(\mathbb Z_5)$.
Moreover every logarithm from $(Q_5^{[5]})'$ lies in
$H_{\mathbb Q}$, so its group commutators have logarithms
in the displayed saturated lattice.

For the reverse containment the program forms actual
group commutators between
\[
 J(b)b^{-1},\qquad J(c_0)c_0^{-1},\qquad
          g_w\quad(2\leq|w|\leq4),
\]
retaining pairs whose weights sum to at most five.
There are $308$ such commutators. Their logarithms are
$5$-integral, lie in $H_{2,\mathbb Q}$, and have rank $252$
modulo $5$. The logarithm of the second-derived subgroup
is an additive $\mathbb Z_5$-module: that subgroup has
class at most two, and the correction
$\frac12[\log g,\log h]$ can be removed using a square
root of the central commutator. Its logarithm therefore
contains the span of the $308$ vectors, which is the
whole lattice by the unit-rank certificate. This proves
the integral equality used to reduce the equations.

Write $X,Y,Z$ for conjugation by $x,y,xy$ and set
\[
 \mathcal N_J(g)=gJ(g)J^2(g),\qquad
 \epsilon(c)=c\mathcal N_J(c)^{-1/3},\qquad c=[x,y].
\]
Here $J$ has order three, as in A.23. On the metabelian
quotient, $\epsilon(c)$ is exactly the nontrivial
$C_3$-character component of $c$.
The four group words
\[
\begin{gathered}
 \mathcal N_J(f),\qquad \mathcal N_J(h),\\
 X(Y(f))X(f)^{-1}Y(f)^{-1}f\,\epsilon(c)^2,\\
 Z(h)h^{-1}\bigl(X(Y(f))X(f)^{-1}\bigr)^{-1}
\end{gathered}                                                   \tag{A--X}
\]
must belong to $(Q_5^{[5]})''$. Together with
\[
 f\Phi(f)=1,\qquad h\Phi(h)=1,\qquad
          xf^{-1}yf=h^{-1}xyh,                                  \tag{A--XI}
\]
they are precisely the seven conditions of A.24, now
expressed in actual group coordinates.
After taking logarithms, project the last four residuals
modulo the saturated lattice just proved, and retain
the $204$ degree-five coordinates in each residual.
All lower-degree residuals vanish by A.24.

All these residuals are affine in (A--IX).
The old corrections have weight four, so their first
quadratic substitution term has weight at least seven;
the new corrections are central of weight five.
For completeness the central derivatives are explicit.
For $w\in L_5$, put $j=J(w)$, $j_2=J^2(w)$, $p=P_1(w)$,
where $P_1$ is induced by the degree-one action in A.23,
and put
\[
 \ell_u=\tfrac23(j+j_2-2w),\qquad
 \ell_v=\tfrac23(2j_2-j-w).
\]
In the order consisting of the two norms, the third
peripheral residual, and the four words (A--X), a change
of the $f$ coordinate by $w$ has derivatives
\[
 (w+p+\ell_u,\ \ell_v,\ j_2-j,\
   w+j+j_2,\ 0,\ j_2-2j+w,\ j-j_2),                           \tag{A--XII}
\]
and a change of the $h$ coordinate by $w$ has derivatives
\[
 (0,\ w+p,\ w-j_2,\ 0,\ w+j+j_2,\ 0,\ j_2-w).              \tag{A--XIII}
\]
To derive the only substitution terms in these formulas,
changing $f$ by $\exp w$ changes the Schreier images of
$B,C,D$ by $(J-1)w,(J^2-J)w,(1-J^2)w$, respectively.
Applying these changes to $u_1,v_1$ gives
$\ell_u,\ell_v$. All other entries follow by multiplying
central corrections and using the deck images of $x,y,xy$.
The four module entries are then projected modulo $H_2$.

The program evaluates all $20$ prior-parameter columns
by exact group formulas and constructs all $408$ central
columns using (A--XII)--(A--XIII). Four declared Lyndon
basis indices, for each of $f$ and $h$, are independently
checked against full group evaluation; these eight
cross-checks support the derived formulas and are not
described as an exhaustive recalculation of all columns.
Every matrix and right-hand-side entry is $5$-integral.

The complete system, particular solution, and all $90$
kernel vectors are saved by
\texttt{a\_metabelian\_lift\_class5.py}.
The independent verifier \texttt{a\_class5\_verify.py}
does not import the series calculation. It verifies
the system hash, $Aq=b$, all $Av_i=0$, the identity
minor on the free coordinates, and rank $338$ modulo
$5$. The independent vectors give the matching rational
upper rank bound. All their entries are $5$-integral,
so the free-coordinate identity also proves that they
give the entire integral kernel. The nonzero Smith
factors are units.

For an explicit description of the projection, let
$q=(q_0,\ldots,q_{19})$ denote the A.24 free coordinates
in their saved order. Its image is
$q=q^{(0)}+s v^{(1)}+t v^{(2)}$, $s,t\in\mathbb Z_5$,
where the following table gives every coordinate.
\[
\begin{array}{r|r|r|r}
i&q_i^{(0)}&v_i^{(1)}&v_i^{(2)}\\\hline
0&5413/2916&1&0\\
1&2041/13122&1/9&5/3\\
2&-2720/2187&-2/3&0\\
3&-61/2187&0&1\\
4&9475/26244&2/9&4/3\\
5&-11317/26244&-2/9&2/3\\
6&-914/6561&-1/9&-8/3\\
7&-3505/13122&-1/9&7/3\\
8&7903/4374&1&1\\
9&5467/8748&1/3&0\\
10&7535/4374&1&5\\
11&-575/972&-1/3&-1\\
12&193/324&1/3&1\\
13&-2791/13122&-1/9&1/3\\
14&1366/729&1&0\\
15&167/243&1/3&-2\\
16&287/486&1/3&1\\
17&-859/1458&-1/3&-1\\
18&-2705/2187&-2/3&0\\
19&-64/2187&0&1
\end{array}                                                      \tag{A--XIV}
\]
The independent verifier computes this image from all
$90$ kernel vectors. Its rank is $2$ over both fields,
and coordinates $0,3$ give a unit identity minor.
It saves the equivalent $18$ affine equations as well.
The kernel consequently has rank $88$.
Substitution of $q=0$ fails the saved affine equations,
so the previously selected class-four particular point
does not extend, although the full family does.

Reduction of the integral parameter equations modulo
$5^n$ gives $5^{90n}$ solutions, $5^{2n}$ extendible
class-four parameter points, and $5^{88n}$ points in
each fibre. These are counts of residue solutions of
the parameter equations. They do not assert that naive
Hall-coordinate reduction is a group quotient at class
five. The group identities over $\mathbb Z_5$ themselves
are exact. The norms imply $\Phi^2=1$, and the preserved
order-$75$ quotient of A.22 proves exact actual and
outer order two. No group elements or parameter
solutions were sampled.
\end{proof}

\begin{statement}{A.26}{Proved}{An all-class relative pro-5 and fibre-product lift}
Let $F=\widehat F(x,y)$,
$K=\ker(F\to C_3)$ with $x,y\mapsto1$, and put
\[
 N=K^{(5)},\qquad R=\ker(K\to N),\qquad Q=F/R.
\]
The image of the specified A.22 metabelian involution
on $Q/Q''$ lifts to an actual continuous automorphism
$\Phi$ of $Q$ of order two. Its outer class also has
order two. There are $f,h\in Q'$ with the prescribed
A.22 images such that
\[
 \Phi(x)=x,\quad\Phi(y)=f^{-1}yf,\quad
 \Phi(xy)=h^{-1}xyh,\quad f\Phi(f)=h\Phi(h)=1.
                                                               \tag{A--XV}
\]
Thus this is a lift through every nilpotency class of
the same open free pro-$5$ kernel, preserving each
oriented peripheral class individually.

More strongly, put $M=F/F''$. The full A.22 action
on $M$, not only its relative pro-$5$ image, lifts to
\[
 \mathcal Q=F/(F''\cap R)
       \cong M\mathbin{\times}_{Q/Q''}Q.                         \tag{A--XVI}
\]
On $\mathcal Q$ it again has exact actual and outer
order two and satisfies (A--XV) with derived
conjugators. This does not assert a lift to $F$, and
does not assert that the all-class lift extends the
particular nonmetabelian pair selected in A.23.
\end{statement}
\begin{proof}
The free pro-$5$ group $N$ has basis $a,b,c_0,d$ as
above. Set $z=a^{-1/3}x$ and $J=\operatorname{Ad}(z)$.
Then $z$ and $J$ have order three, $J(a)=a$, and
$Q=N\rtimes\langle z\rangle$.
The three peripheral cubes in $N$ can be taken as
\[
       a=x^3,\qquad b_0=y^3=bc_0d,\qquad
       c_\infty=(xy)^{-3}=(c_0abd)^{-1}.                        \tag{A--XVII}
\]
All three are primitive. For $b_0$ one can replace $d$
in the free basis; for $c_\infty^{-1}$ one can replace
$a$. These are elementary Nielsen changes of basis.

\emph{The precise boundary-preserving subgroup.}
Let $\mathcal C\subset\operatorname{Aut}(N)$ consist
of the automorphisms fixing $a$ exactly and preserving
the individual $N$-conjugacy classes of $b_0,c_\infty$,
with their displayed orientations. This is a closed
profinite subgroup: conjugacy classes in a compact
group are compact, and evaluation is continuous.
It contains $J$. One can see the conjugacy assertion
from the degree-three cyclic covering of the tripod:
each cusp has a unique lift. Algebraically the
order-three element $b_0^{-1/3}y$ centralizes $b_0$
and has the same deck image as $z$, so conjugation
by $z$ takes $b_0$ to an $N$-conjugate. For $c_\infty$
use the square of the order-three part of $xy$.

Write $V=N^{\mathrm{ab}}=\mathbb Z_5^4$, with letters
$A,B,C,D$, and let
\[
 W=\langle A,B+C+D\rangle,\qquad
 a_*=A,\quad b_*=B+C+D,\quad u=B,\quad v=C.
\]
The discrete surface underlying this cover has genus
one and three boundary components: its Euler
characteristic is $-3$, and each of the three base
boundaries has one lift. The radical of its
intersection form is precisely $W$, while $V/W$
is the rank-two proper-homology module.
Define the integral $5$-adic involution $P$ by
\[
 P(a_*)=a_*,\quad P(b_*)=b_*,\qquad
 P(u)=-u+\tfrac23(b_*-a_*),\quad
 P(v)=-v+\tfrac23(b_*-a_*).                                   \tag{A--XVIII}
\]
It is the degree-one action used in A.23. Directly,
$J(u)=v$ and $J(v)=b_*-a_*-u-v$, so $P$ commutes
with $J$. Their reductions generate $C_6$ in
$\operatorname{GL}(V/5V)$.

Here is a direct proof that $\overline P$ lies in the
image of $\mathcal C$. Pure mapping classes fixing
all three boundaries pointwise act on the discrete
fundamental group by automorphisms fixing the based
first boundary loop exactly and preserving the other
two conjugacy classes. They therefore give elements
of $\mathcal C$ after pro-$5$ completion. Two handle
twists generate $\operatorname{SL}_2(\mathbb Z)$
on proper homology. The required additional shears
are also explicit. If $r$ is either of the other
two boundary classes, band a simple handle curve
of class $w$ to that boundary to obtain a simple
curve of class $w+r$. The twist transvection
$t_w(s)=s+\langle s,w\rangle w$, obtained by counting
signed crossings through its annulus, gives
\[
 t_{w+r}t_w^{-1}(s)=s+\langle s,w\rangle r.                    \tag{A--XIX}
\]
Using two handle classes and the two independent
boundary classes gives every map in
$\operatorname{Hom}(V/W,W)$ integrally. Thus these
mapping classes realize every integral homology
automorphism fixing $W$ pointwise and preserving
the proper intersection form. In (A--XVIII), replace
$2/3$ by the integer $4$; its reduction modulo $5$
is $\overline P$. This proves the required inclusion
without asserting that a closed-surface symplectic
surjectivity theorem already proves the boundary
claim. For comparison, the closed-surface theorem
and its proof are given by
\citet[Theorem 0.1]{A_PutmanSymplectic}; the radical
calculation needed here is (A--XIX).

\emph{Coprime lifting with the deck action fixed.}
Let $\rho:\mathcal C\to\operatorname{GL}(V/5V)$.
Its kernel is pro-$5$, since it is a closed subgroup
of the congruence kernel of the automorphism group
of a finitely generated pro-$5$ group. This precise
fact is recalled in
\citet[\S2, preceding Definition 2.1]{A_HajirMaire2016}.
The inverse image of
$\langle\overline J,\overline P\rangle\cong C_6$
therefore has a complement by the profinite
Schur--Zassenhaus theorem
\citep[Theorem 2.5]{A_HajirMaire2016}.
Its order-three subgroup and $\langle J\rangle$
are complements in $\rho^{-1}(\langle\overline J\rangle)$.
Conjugacy of complements by $\ker\rho$ lets us
choose the order-six complement to contain exactly
$J$. The projection distinguishes $J$ from $J^2$.
The order-two element $\Phi_N$ of this cyclic
complement commutes with $J$, belongs to
$\mathcal C$, and induces $\overline P$.
The finite $C_6$ constructed here is a subgroup
of $\operatorname{Aut}(N)$ for the pro-$5$ group
$N$. The discrete mapping classes were used only
to realize its prescribed modulo-$5$ homology matrices.
Extend it to $Q$ by $\Phi(z)=z$.
Because it fixes $a$, it fixes $x=a^{1/3}z$.

In fact its action on $V$ is exactly (A--XVIII).
The invariant part $V^J=W$ is primitive, and $3$
is a $5$-adic unit, so $V=W\oplus(1-(1+J+J^2)/3)V$.
It acts as the identity on $W$ and commutes with $J$.
On the complementary rank-two part it has square
one and reduction $-1$. The idempotent $(\Phi_N+1)/2$
on that summand has zero reduction; Nakayama's lemma
therefore makes its image zero. This proves exact
action $-1$ on that summand and hence exact action $P$.

\emph{From boundary cubes to boundary elements.}
For a primitive element $b_0$ in a free pro-$5$
group, $C_N(b_0)=\langle b_0\rangle$. The Magnus
centralizer argument used in A.19 proves this in
rank four as well: with $b_0$ as one basis element,
a series commuting with its variable involves only
that variable; the retraction to its procyclic
subgroup has the same series, so injectivity puts
the original element in that subgroup.
It follows that $C_Q(y^3)=\langle y\rangle$:
multiply a centralizing element by a power of $y$
to put it in $N$, then use this primitive centralizer.
The cube equation together with the specified
nontrivial $C_3$ coset has exactly one root in this
centralizer. Cube roots in $Q$ without a specified
coset need not be unique.
Consequently preservation of the $N$-conjugacy
class of $y^3$ implies preservation of the
$N$-conjugacy class of $y$. The identical argument
applies to $(xy)^3$, whose deck image is also
nontrivial. Hence some $f,h\in N$ satisfy the
first three equations (A--XV).

\emph{Derived conjugators and exact norms.}
The abelianization of $Q$ is
\[
                 A_Q=Q^{\mathrm{ab}}\cong\mathbb Z_5^2\times C_3.
                                                               \tag{A--XX}
\]
Indeed the $J$-coinvariants of $V$ have basis $A,B$,
with $C=B,D=A+B$. The images of $x^3,y^3$ are
$A,A+3B$, a basis because $3$ is a unit.
Since $\Phi$ fixes $x^3$, preserves $y^3$ up to
conjugacy, and fixes $z$, it is the identity on
$A_Q$. Its square is the identity on $Q$.
Thus $f\Phi(f)$ centralizes $y$ and lies in $N$,
so $f\Phi(f)=(y^3)^t$ for some $t\in\mathbb Z_5$.
Replace $f$ by $(y^3)^{-t/2}f$. Its conjugation
of $y$ is unchanged, and its new norm is one:
use $\Phi(y^3)=f^{-1}y^3f$ in the product.
In $A_Q$ the old norm equation was
$2[f]=t[y^3]$; the replacement therefore has
zero abelianization, as required.
Applying the same argument to $h$, with $(xy)^3$
in place of $y^3$, gives $h\in Q'$ and norm one.
These replacements do not change the actual
involution $\Phi$.

\emph{The exact relative Alexander module.}
Put $\Lambda=\mathbb Z_5[[A_Q]]$. The full profinite
Magnus theorem used in A.22 gives a cyclic
surjection
\[
       \Lambda c\longrightarrow Q'/Q'',\qquad c=[x,y].          \tag{A--XXI}
\]
One can also obtain cyclicity directly from the
normal generation of $F'$ by $[x,y]$; the coefficient
action factors through $A_Q$, and the target is
pro-$5$. To prove that (A--XXI) has no extra
relations, specialize the Magnus representation
to $\Lambda^2\rtimes A_Q$, sending
$x$ to $(e_x,X)$ and $y$ to $(e_y,Y)$.
The kernel of this semidirect product over $C_3$
is pro-$5$, so the closed subgroup generated by
these two images is a genuine quotient of $Q$.
The commutator has additive vector
\[
                    (1-Y)e_x+(X-1)e_y.                          \tag{A--XXII}
\]
For the full profinite formula see
\citet[Propositions 3.2.1 and 4.1.1(b), Corollary 4.1.5]{A_ChenDeligne2017}.
The specialization here proves the relative
statement without assuming that completion
preserves a discrete Alexander computation.

Let $U=X^3,V_0=Y^3$. The two ring factors are
\[
 \Lambda\cong
 \mathbb Z_5[[U-1,V_0-1]]\ \times\
 \mathbb Z_5[\zeta_3][[U-1,V_0-1]].                            \tag{A--XXIII}
\]
In the first factor $X=U^{1/3},Y=V_0^{1/3}$;
in the second they are
$X=\zeta_3U^{1/3},Y=\zeta_3V_0^{1/3}$.
Both factors are domains. The element $X-1$
is nonzero in the first and is a unit in the
second. Thus the vector (A--XXII) has zero
annihilator. Composing (A--XXI) with this Magnus
map proves
\[
                         Q'/Q''=\Lambda c.                     \tag{A--XXIV}
\]
The idempotent $e$ selecting the second factor is
exactly the image of the A.22 Haar idempotent:
the surviving prime-to-$5$ characters satisfy
$X=Y=\zeta_3$ or $\zeta_3^2$.

The normalized $f$ has image $c^r$ for a unique
$r\in\Lambda$. Since $\Phi$ is the identity on
$A_Q$, its action on (A--XXIV) is multiplication by
\[
                 \theta=1+(X-1)(Y-1)r.                         \tag{A--XXV}
\]
This follows by the same commutator calculation
as A.22. It satisfies $\theta^2=1$. Each of the
two domain factors therefore has sign $+1$ or
$-1$. In the first factor the augmentation of
(A--XXV) is $1$, so its sign is $+1$.
In the second factor, set $5=U-1=V_0-1=0$.
The class of $c$ in $N^{\mathrm{ab}}/5$ is
$C-B$, which generates its nontrivial $C_3$
character over $\mathbb F_{25}$. The action $P$
is $-1$ there, forcing the second sign to be
$-1$. Hence $\theta=1-2e$ exactly.
Because $(X-1)(Y-1)$ is a nonzerodivisor and
is a unit on the $e$-factor,
\[
       r=\frac{-2e}{(X-1)(Y-1)}.                               \tag{A--XXVI}
\]
The third peripheral equation gives, for the
unique residue $h=c^s$,
\[
       (XY-1)s=X(Y-1)r,
       \qquad s=\frac{X(Y-1)r}{XY-1}.                          \tag{A--XXVII}
\]
Here $XY-1$ is a nonzerodivisor on the first
factor and a unit on the second, so $s$ is
also unique. Equations (A--XXVI)--(A--XXVII)
are precisely the specified A.22 residues.

The involution is outer already on $N^{\mathrm{ab}}$.
An inner automorphism of $Q$ acts there by a
power of $J$, whereas $\Phi$ is $-1$ on the
nontrivial character summand and
$-1\notin\langle\zeta_3\rangle$.
Equivalently its preserved order-$75$ quotient
is the noninner automorphism verified in A.22.
Thus both its actual and outer orders are two.

\emph{Lifting the whole metabelian action.}
The common quotient of $M=F/F''$ and $Q=F/R$
is $F/(F''R)=Q/Q''$; the product $F''R$ is
closed because both factors are compact normal
subgroups. The elementary normal-subgroup
pullback lemma gives (A--XVI): a compatible
pair of quotient elements has lifts differing
by an element of $F''R$, which can be adjusted
by its two factors to give one simultaneous
lift, and the kernel is $F''\cap R$.
The A.22 automorphism on $M$ and the constructed
$\Phi$ on $Q$ induce the same action on this
common quotient by (A--XXVI)--(A--XXVII).
They therefore act componentwise on the
fibre product. Their respective normalized
$f$ and $h$ also have identical common images,
so give conjugators in the fibre product
satisfying both norm equations exactly.
They lie in its derived subgroup: its
abelianization maps isomorphically to
$M^{\mathrm{ab}}=F^{\mathrm{ab}}$, since
$F''\cap R\subset F''$.
The order-$75$ quotient again proves exact
outer order two.
\end{proof}

\paragraph{The remaining full profinite obligation.}
The all-class result concerns the relative pro-$5$
quotient and the fibre product (A--XVI). Both are
prosolvable; the free profinite group also has
nonabelian simple finite quotients, so this
construction has not recovered the full group.
The concrete residual in A.12 is whether there
are compatible peripheral order-two lifts of
the A.22 action through the remaining finite
quotients, equivalently a corresponding actual
automorphism of $\widehat F_2$ under the
normalization of A.19. No obstruction can be
deduced merely from increasing the nilpotency
class in this particular open pro-$5$ kernel:
A.26 supplies one compatible lift through all
of them. The extension of the specifically
chosen A.23 pair to all classes is a separate
problem; A.25 proves it through class five.
